\section{Conclusion}
\label{sec:conclusion}

We presented the Wiggle Framework, a three-axis diagnostic for measuring LLM-as-judge reliability, and validated it empirically across 9 frontier models on 384 borderline safety items.  Five experiments revealed that the axes are largely independent, that conviction degradation curves expose qualitatively distinct failure archetypes (Cliff, Instant Capitulator, Immovable Object, Gradual Yielder), and that the correct-vs-wrong-label persistence test successfully disambiguates sycophancy from rational updating---finding that seven of nine models are rational.

Three findings have direct implications for safety evaluation practice.  First, the universal restrictive bias means that any challenge-based review protocol will systematically inflate unsafe counts; prevalence estimates should account for this directional asymmetry.  Second, stubbornness is not reliability: the most persistent model (Grok-4.1) is the only truly sycophantic one, while the least persistent (Claude~Opus) is the most informationally responsive.  Third, eval-awareness modulates judge behavior, meaning benchmark-calibrated thresholds will not transfer directly to naturalistic deployment.

As LLM-as-judge systems become load-bearing infrastructure for safety evaluation, content moderation, and model launch decisions, the question is no longer \emph{how accurate is this judge?} but \emph{how much does this judge wiggle, and in what ways?}  The Wiggle Framework equips practitioners with the diagnostic toolkit to answer that question.
